\section{Poisson summation and the dual mean square}\label{sec:initialization}

We now turn to the mean-square estimate in Proposition~\ref{thm:ms}.
Following Step~3 of the outline, we expand the square and apply
Poisson summation in the row variable $u$. The finite Fourier
transforms of the characters supply Gauss sums. Arithmetic identities
then combine these Gauss sums with the original M\"obius coefficients
to give the normalized cubic Gauss-sum coefficients that appear in the theta
function.

The column indices produced by expanding the square need not be
coprime. We first extract their common factor, then use M\"obius
inversion to separate the remaining coprimality condition. These
operations introduce an auxiliary twisting index and an exclusion
ideal. The comparison below removes the exclusion without changing
the row range or the product of the column and auxiliary scales.

We record the arithmetic and Poisson identities first, then define
this family. Proposition~\ref{prop:poisson-reduction} states the
estimate for it that suffices to prove \eqref{eq:ms}; the rest of
the section proves that implication.

\subsection{The arithmetic identities}
We first record the identities that convert between M\"obius
coefficients and normalized cubic Gauss sums. Recall that, for a
squarefree primary element $n$ with $(n,S)=1$,
\[
 \alpha(n)=\frac{n}{|n|},\qquad
 \gamma_j(n)=\frac{1}{\sqrt{\NK(n)}}
       \sum_{x\bmod n}\chi_n(x)^j e(x/n)\quad(j\in\mathbb Z),
\]
where $e(z)=\exp(4\pi\mathrm i\,\operatorname{Im}z/\sqrt3)$, as in
\eqref{eq:intro-gauss-sum}. Put
\[
 a_\xi(n)=\overline{\alpha(n)}\gamma_2(n)\xi(n)
 \qquad(n\text{ squarefree}),
\]
where $\xi$ is a fixed ray class character. The character $\xi$ accounts for the
residue-class factors in the identities below.

By character orthogonality \cite[\S3.1]{IK04}, we can absorb functions
on a fixed ray class group into a finite sum of twists $\xi$.
Choose a fixed modulus, supported on $S$, divisible by $12$
and sufficiently large that all reciprocity factors below
depend only on the corresponding ray classes. We then expand these
factors in characters of this finite group.

The following lemma collects standard consequences of the Gauss--Jacobi
identities, quadratic Gauss-sum evaluations, and reciprocity.

\begin{lemma}\label{lem:arithmetic}
On a fixed ray class group, there are a function $G$ with values in
$\{z\in\mathbb C:|z|=1\}$ and a symmetric $\{\pm1\}$-valued
bicharacter $\mathcal R$. This means that $\mathcal R(a,b)=\mathcal R(b,a)$ and $\mathcal R$ is
multiplicative in each argument separately:
\[
 \begin{aligned}
 \mathcal R(aa',b)&=\mathcal R(a,b)\mathcal R(a',b),\\
 \mathcal R(a,bb')&=\mathcal R(a,b)\mathcal R(a,b'),
 \end{aligned}
\]
for all classes $a,a',b,b'$ in the ray class group.
These functions have the following properties.
Let $a,b,n$ be primary elements prime to $S$, with $(a,b)=1$ and
$n$ squarefree. Then
\begin{align}
 \chi_b(a)&=\mathcal R(a,b)\chi_a(b),&
 G(ab)&=G(a)G(b)\mathcal R(a,b),                                      \label{eq:recip}\\
 \gamma_2(n)^3&=\mu(n)\alpha(n),&
 \gamma_1(n)\gamma_2(n)&=\mu(n)\alpha(n)G(n),                 \label{eq:gj}\\
 G(n)&=\overline{\chi_n(4)}\gamma_3(n),&
 \gamma_1(n)\gamma_{-1}(n)&=\chi_n(-1).                     \label{eq:g}
\end{align}
Consequently
\begin{align}
 \overline{\alpha(n)}\gamma_2(n)\gamma_1(n)&=\mu(n)G(n),\label{eq:convert1}\\
 \mu(n)\gamma_{-1}(n)
 &=\chi_n(-1)G(n)^{-1}\overline{\alpha(n)}\gamma_2(n),\label{eq:convert2}\\
 a_\xi(ab)&=a_\xi(a)a_\xi(b)\chi_b(a)^4,                     \label{eq:crt-a}\\
 \chi_a(-1)\overline{G(a)}G(b)\mathcal R(a,b)&=G(ba^{-1}).            \label{eq:quotient}
\end{align}
In \eqref{eq:crt-a}, $a,b$ are also squarefree.
The identity for $G(ab)$ in \eqref{eq:recip} extends to all classes of the fixed ray class
group; the quotient in \eqref{eq:quotient} is taken in that group.
\end{lemma}
The proof, including the dependence of $G$ and $\mathcal R$ on fixed ray
classes, is given in Appendix~\ref{app:gauss-identities}.
Under Poisson summation, \eqref{eq:convert2} converts the M\"obius
coefficients to $a_\xi(n)$, up to fixed ray class factors, while
\eqref{eq:convert1} converts them back.

\subsection{Poisson summation with excluded primes}
The row sum to which we apply Poisson summation will have an
additional coprimality restriction. We record the formula with that
restriction included, so that its effect on the Fourier frequencies
and the normalization is explicit.

For a nonzero ideal $\mathfrak r$, write $\rad \mathfrak r$ for the product of
its distinct prime divisors. We use the additive character $e$
introduced in Step~3. In quotients and Gauss sums, use a fixed
generator for each ideal, chosen primary when the ideal is prime to $3$.

\begin{lemma}\label{lem:poisson}
Let $\mathcal H>0$, let $\mathfrak r$ be a nonzero ideal, and let
$\Phi(\NK(k)/\mathcal H)$ be a smooth radial Schwartz weight
on the row lattice. Let $\chi$ be a primitive multiplicative character
of $(\OO/\mathfrak m)^\times$, viewed as a function on $\OO$ by
reduction modulo $\mathfrak m$ and extension by zero on nonunits.
Write 
\[
\gamma(\chi)=\NK(\mathfrak m)^{-1/2}
\sum_{x\bmod\mathfrak m}\chi(x)e(x/\mathfrak m)
\]
for its normalized Gauss sum.
Since $\chi$ is primitive, its modulus $\mathfrak m$ is determined by $\chi$ and is suppressed in the notation. Then
\begin{equation}\label{eq:poisson}
 \begin{aligned}
 &\sum_k\chi(k)\ind_{(k,\mathfrak r)=1}\Phi(\NK(k)/\mathcal{H})\\
 &\quad=\frac{\mathcal{H}\gamma(\chi)}{\sqrt{\NK(\mathfrak m)}}
   \sum_{d\mid\rad \mathfrak r}\frac{\mu(d)\chi(d)}{\NK(d)}
   \sum_h\overline{\chi(h)}
        \widehat\Phi\!\Bigl(\frac{\mathcal{H}\NK(h)}{\NK(d)\NK(\mathfrak m)}\Bigr).
 \end{aligned}
\end{equation}
Here $\widehat\Phi$ is defined by
\[
 \widehat\Phi(|w|^2)=\frac2{\sqrt3}\int_{\mathbb C}
       \Phi(|z|^2)e(-zw)\,dx\,dy,\qquad z=x+\mathrm i y.
\]
The measure is normalized so that $\OO$ has covolume one.
Here $k,h\in\OO$, and $h$ is the Fourier frequency. We have
\[
 \chi\text{ nonprincipal}\quad\Longrightarrow\quad
 \overline{\chi(0)}=0.
\]
For the principal primitive character ($\chi=\mathbf1$, $\mathfrak m=1$),
the zero-frequency contribution is
\[
 \mathcal H\widehat\Phi(0)
   \prod_{\mathfrak p\mid \mathfrak r}\Bigl(1-\frac1{\NK(\mathfrak p)}\Bigr).
\]
\end{lemma}

\begin{proof}
Inclusion--exclusion followed by $k=d\ell$ gives
\[
 \sum_k\chi(k)\ind_{(k,\mathfrak r)=1}\Phi(\NK(k)/\mathcal{H})
 =\sum_{d\mid\rad \mathfrak r}\mu(d)\chi(d)
       \sum_\ell\chi(\ell)\Phi\!\Bigl(\frac{\NK(\ell)}{\mathcal{H}/\NK(d)}\Bigr).
\]
Apply lattice Poisson summation \cite[Theorem~4.5]{IK04} in $\ell$
on residue classes modulo $\mathfrak m$, at scale $\mathcal H/\NK(d)$.
The primitive Gauss-sum identity \cite[(3.12)]{IK04} evaluates the
finite Fourier transform as
$\sqrt{\NK(\mathfrak m)}\gamma(\chi)\overline{\chi(h)}$
for every $h$, giving \eqref{eq:poisson}. For a nonprincipal primitive
character, $\overline{\chi(0)}=0$, so the zero-frequency term vanishes.
For the principal character, sum $\mu(d)/\NK(d)$ over
$d\mid\rad \mathfrak r$ to obtain the displayed zero-frequency contribution.
\end{proof}

\subsection{The dual mean squares} 
We use the following mean square of the column sums with
coefficients $a_\xi(n)$.

\begin{definition}
The parameters $\mathcal H$, $X$, and $F$ are the norm scales of
$k$, $n$, and $f$, respectively. Here $k$ ranges over elements of
$\OO$, and a star restricts a sum to squarefree ideals prime to $S$,
represented by their primary generators. For a smooth compactly
supported weight $W$ on $(0,\infty)$, define the \emph{dual mean square} by
\begin{equation}\label{eq:energy}
 \begin{aligned}
 &\E(\mathcal{H},X,F;\xi,W)\\
 &\qquad=\frac1{XF}
 \sum_{\substack{F\le\NK(f)<2F\\(f,S)=1}}^*\sum_{0<\NK(k)\le \mathcal{H}}
 \Bigl|\sum_{(n,S)=1}^*a_\xi(n)\chi_n(k)\chi_n(f)^4
                  W(\NK(n)/X)\Bigr|^2.
 \end{aligned}
\end{equation}
\end{definition}

Below, $\xi$ ranges over all characters of the fixed ray class group
chosen above, through which $G$ and $\mathcal R$ factor.

Write $\E_{\mathfrak r}$ for the same expression with the additional
restriction $(n,\mathfrak r)=1$. This notation is only needed in the
Poisson reductions; the following comparison returns to $\E$.

\begin{lemma}\label{lem:remove-exclusions}
Let $r_0$ be the product of the primes dividing $\mathfrak r$ outside
$S$. For $\mathcal H,X>0$, $F\ge1$, and smooth compactly supported $W$,
\begin{equation}\label{eq:remove-exclusions}
 \E_{\mathfrak r}(\mathcal H,X,F;\xi,W)
 \le\tau_{\mathrm{div}}(r_0)
 \sum_{d\mid r_0}\E\!\Bigl(\mathcal H,\frac X{\NK(d)},F\NK(d);\xi,W\Bigr).
\end{equation}
\end{lemma}
\begin{proof}
Let $C_{\mathfrak r}(X;k,f)$ denote the inner column sum defining
$\E_{\mathfrak r}$, with $\xi,W$ fixed. Inclusion--exclusion,
$n=dm$, and \eqref{eq:crt-a} give
\[
 \begin{aligned}
 C_{\mathfrak r}(X;k,f)
 &=\sum_{d\mid r_0}\mu(d)
   \sum_{\substack{(n,S)=1\\d\mid n}}^*
   a_\xi(n)\chi_n(k)\chi_n(f)^4W(\NK(n)/X)\\
 &=\sum_{\substack{d\mid r_0\\(d,f)=1}}
   \mu(d)a_\xi(d)\chi_d(k)\chi_d(f)^4
   C_1(X/\NK(d);k,df).
 \end{aligned}
\]
Indeed, $\chi_m(d)^4$ enforces $(m,d)=1$ and combines with
$\chi_m(f)^4$; terms with $(d,f)\ne1$ vanish. Each exterior
coefficient has modulus at most one. Apply Cauchy--Schwarz in $d$,
then enlarge the injective image $f\mapsto df$ to the squarefree
range $F\NK(d)\le\NK(df)<2F\NK(d)$. The normalizing product is
unchanged: $(X/\NK(d))(F\NK(d))=XF$.
\end{proof}


\subsection{Reduction to the dual mean square}
The following proposition gives the dual estimate sufficient for
\eqref{eq:ms}. Its proof applies Poisson summation directly to the
original M\"obius sums.

\begin{proposition}\label{prop:poisson-reduction}
Let $\vartheta=1/10$ and put $H=D^{1+\vartheta}=D^{11/10}$.
For each fixed $C\ge1$ and all real $B,F\ge1$, consider the ranges
\begin{equation}\label{eq:initial-scales}
 X=\frac D{BF},\qquad 0<\mathcal{H}\le\frac{CD^2}{HB^2},\qquad
 XF=\frac DB.
\end{equation}
Suppose that for every $\varepsilon>0$ there exists an integer
$J=J(\varepsilon)\ge1$ such that, for every compact interval
$I\subset(0,\infty)$ and every smooth $W$ supported in $I$,
\begin{equation}\label{eq:auxiliary-target}
 \E(\mathcal H,X,F;\xi,W)
 \ll_{I,C,\varepsilon}\|W\|_{C^J(I)}^2D^\varepsilon XF,
\end{equation}
where
\[
 \|W\|_{C^J(I)}=\max_{0\le j\le J}\sup_{x\in I}|W^{(j)}(x)|.
\]
For functions of several variables, the norm uses all partial derivatives
of total order at most $J$.
Then the original mean-square
estimate in Proposition~\ref{thm:ms} holds.
\end{proposition}

\begin{proof}
Write $I=[a_0,b_0]$. Choose a nonnegative radial Schwartz weight
$\Phi$ such that $\Phi(x)\ge1$ on $[0,1]$ and
$\operatorname{supp}\widehat\Phi\subset[0,C_\Phi]$ for some fixed $C_\Phi>0$.
It suffices to prove
$\mathcal M_D\ll_{I,\varepsilon}
\|W\|_{C^{J'}(I)}^2HD^\varepsilon$ for some $J'=J'(\varepsilon)$, where
\begin{equation}\label{eq:expanded-ms}
 \mathcal M_D:=\frac1D\sum_{u\in\OO}\Phi(\NK(u)/H)
 \Bigl|\sum_{(n,S)=1}\mu(n)\chi_n(u)W(\NK(n)/D)\Bigr|^2.
\end{equation}
Conjugate the inner sum, expand the square, and put
$g=(n_1,n_2)$ and $n_j=gz_j$. The squarefree indices satisfy
$(z_1,z_2)=(z_1z_2,g)=1$, and
\[
 \overline{\chi_{n_1}(u)}\chi_{n_2}(u)
 =\ind_{(u,g)=1}\overline{\chi_{z_1}(u)}\chi_{z_2}(u).
\]
Apply Lemma~\ref{lem:poisson} with
$\chi=\overline{\chi_{z_1}}\chi_{z_2}$ and exclusion $g$:
\begin{equation}\label{eq:initial-poisson-rows}
 \begin{aligned}
 &\sum_u\Phi(\NK(u)/H)\ind_{(u,g)=1}\chi(u)\\
 &\quad=\sum_{e\mid g}
 \frac{H\mu(e)\chi(e)\gamma(\chi)}
      {\NK(e)\sqrt{\NK(z_1)\NK(z_2)}}
 \sum_h\overline{\chi(h)}
 \widehat\Phi\!\Bigl(\frac{H\NK(h)}{\NK(e)\NK(z_1)\NK(z_2)}\Bigr).
 \end{aligned}
\end{equation}
The zero frequency occurs only when $z_1=z_2=1$ and contributes
\[
 Z=\frac HD\widehat\Phi(0)
 \sum_{(g,S)=1}^*|W(\NK(g)/D)|^2
       \prod_{p\mid g}\Bigl(1-\frac1{\NK(p)}\Bigr)
 \ll_{\Phi,I}H\|W\|_\infty^2.
\]
Retain all $h\ne0$, including those with $z_1=z_2=1$.

Expand the fixed ray class function
$G(t^{-1})=\sum_\xi c_\xi\xi(t)$.
The Chinese remainder theorem and
\eqref{eq:convert2}--\eqref{eq:quotient} give
\begin{equation}\label{eq:initial-paired-gauss}
 \mu(z_1)\mu(z_2)
       \gamma(\overline{\chi_{z_1}}\chi_{z_2}) =\sum_\xi c_\xi a_\xi(z_1)\overline{a_\xi(z_2)}.
\end{equation}
Write $N=\NK$ for the remainder of this proof, and put
$W_0(x)=x^{-1/2}\overline{W(x)}$. Substituting
\eqref{eq:initial-paired-gauss} into \eqref{eq:initial-poisson-rows},
including the normalization $1/D$ in \eqref{eq:expanded-ms}, gives
\[
 \mathcal M_D-Z=\sum_\xi c_\xi\mathcal S_\xi,
 \qquad
 |\mathcal M_D-Z|\ll\max_\xi|\mathcal S_\xi|,
\]
since the character sum is fixed and finite. Fix $\xi$. Using
$\overline{\chi_z(e)}=\chi_z(e^5)$, its contribution is
\[
 \begin{aligned}
 \mathcal S_\xi
 ={}&\frac H{D^2}
 \sum_g^*\sum_{e\mid g}\frac{\mu(e)N(g)}{N(e)}
 \sum_{h\ne0}
 \sum_{\substack{z_1,z_2\\(z_1,z_2)=1\\(z_1z_2,g)=1}}^*
 a_\xi(z_1)\overline{a_\xi(z_2)}
 \chi_{z_1}(he^5)\overline{\chi_{z_2}(he^5)}
 \\
 &\quad{}\times
 W_0\!\Bigl(\frac{N(gz_1)}D\Bigr)
 \overline{W_0\!\Bigl(\frac{N(gz_2)}D\Bigr)}
 \widehat\Phi\!\Bigl(
 \frac{HN(h)}{N(e)N(z_1)N(z_2)}
 \Bigr).
 \end{aligned}
\]
Here and below every starred variable is squarefree, primary, and
prime to $S$, while $h,k$ range over nonzero elements of $\OO$.
We must show
$|\mathcal S_\xi|\ll_{I,\varepsilon}
\|W\|_{C^{J'}(I)}^2HD^\varepsilon$.

Insert the coprimality identity and change variables:
\[
 \ind_{(z_1,z_2)=1}
 =\sum_{v\mid z_1,\ v\mid z_2}\mu(v),
 \qquad z_j=vm_j.
\]
By \eqref{eq:crt-a}, the common factor from the two columns satisfies
\[
 a_\xi(vm)=a_\xi(v)a_\xi(m)\chi_m(v)^4,
 \qquad
 |a_\xi(v)\chi_v(he^5)|^2=\ind_{(v,h)=1},
\]
where $(v,e)=1$. Thus the same sum becomes
\[
 \begin{aligned}
 \mathcal S_\xi
 ={}&\frac H{D^2}
 \sum_{\substack{g,v\\(g,v)=1}}^*
 \sum_{e\mid g}\frac{\mu(e)\mu(v)N(g)}{N(e)}
 \sum_{\substack{h\ne0\\(h,v)=1}}
 \sum_{\substack{m_1,m_2\\(m_1m_2,gv)=1}}^*
 a_\xi(m_1)\overline{a_\xi(m_2)}
 \\
 &\quad{}\times
 \chi_{m_1}(he^5v^4)\overline{\chi_{m_2}(he^5v^4)}
 W_0\!\Bigl(\frac{N(gvm_1)}D\Bigr)
 \overline{W_0\!\Bigl(\frac{N(gvm_2)}D\Bigr)}
 \\
 &\quad{}\times
 \widehat\Phi\!\Bigl(
 \frac{HN(h)}{N(e)N(v)^2N(m_1)N(m_2)}
 \Bigr).
 \end{aligned}
\]
There is now no restriction $(m_1,m_2)=1$.
Make the bijective change of variables
\[
 \begin{gathered}
 b=g/e,\qquad f=ev,\qquad k=eh,\\
 e=(f,k),\qquad v=f/e,\qquad g=be,\qquad h=k/e.
 \end{gathered}
\]
The resulting $b,f$ are coprime and squarefree, $k\ne0$ is arbitrary,
and
\[
 gv=bf,\qquad he^5v^4=kf^4,\qquad
 \mu(e)\mu(v)=\mu(f),\qquad
 \frac{N(g)}{N(e)}=N(b).
\]
Consequently,
\begin{equation}\label{eq:initial-column-output}
 \begin{aligned}
 \mathcal S_\xi
 ={}&\frac H{D^2}
 \sum_{\substack{b,f\\(b,f)=1}}^*\mu(f)N(b)
 \sum_{k\ne0}
 \sum_{\substack{m_1,m_2\\(m_1m_2,b)=1}}^*
 a_\xi(m_1)\overline{a_\xi(m_2)}
 \chi_{m_1}(kf^4)\overline{\chi_{m_2}(kf^4)}
 \\
 &\quad{}\times
 W_0\!\Bigl(\frac{N(bfm_1)}D\Bigr)
 \overline{W_0\!\Bigl(\frac{N(bfm_2)}D\Bigr)}
 \widehat\Phi\!\Bigl(
 \frac{HN(k)}{N(f)^2N(m_1)N(m_2)}
 \Bigr).
 \end{aligned}
\end{equation}
The characters enforce $(m_j,f)=1$, leaving only the displayed
exclusion $(m_j,b)=1$.
The supports of $W_0$ and $\widehat\Phi$ give
\[
 N(b)N(f)\le b_0D,\qquad
 0<N(k)\le\frac{C_\Phi b_0^2D^2}{HN(b)^2}.
\]

Partition $B\le N(b)<2B$ and $F\le N(f)<2F$ into dyadic ranges.
Denote the corresponding contribution to
\eqref{eq:initial-column-output} by $\mathcal S_{\xi;B,F}$, and put
\[
 X=\frac D{BF},\qquad
 \mathcal H=\frac{C_{I,\Phi}D^2}{HB^2},\qquad
 C_{I,\Phi}=\max(2,C_\Phi b_0^2).
\]
These satisfy \eqref{eq:initial-scales}.
In the variables $x_j=N(m_j)/X$, the coupled smooth weight is
\[
 \begin{gathered}
 \mathcal K_{b,f,k}(x_1,x_2)
 =W_0(r x_1)\overline{W_0(r x_2)}
 \widehat\Phi\!\Bigl(\frac{a}{x_1x_2}\Bigr),\\
 r=\frac{N(b)N(f)}{BF}\in[1,4),\qquad
 a=\frac{HN(k)}{N(f)^2X^2}\le C_{I,\Phi}.
 \end{gathered}
\]
For every integer $q\ge0$, these kernels satisfy
\[
 \sup_{b,f,k}\|\mathcal K_{b,f,k}\|_{C^q([a_0/4,b_0]^2)}
 \ll_{I,\Phi,q}\|W\|_{C^q(I)}^2.
\]
For $U\in C_c^\infty(I_*)$, where $I_*=[a_0/8,2b_0]$, put
\[
 S_U(b,f,k)=\sum_{\substack{(n,S)=1\\(n,b)=1}}^*a_\xi(n)
       \chi_n(k)\chi_n(f)^4U(N(n)/X).
\]
The shifts $(X,F)\mapsto(X/N(d),FN(d))$ preserve
\eqref{eq:initial-scales}. Thus Lemma~\ref{lem:remove-exclusions}
and \eqref{eq:auxiliary-target}, with $J=J(\varepsilon/4)$, give
\[
 \begin{aligned}
 \sum_{f,k}|S_U(b,f,k)|^2
 &=XF\,\E_{(b)}(\mathcal H,X,F;\xi,U)\\
 &\le XF\,\tau_{\mathrm{div}}(b)
       \sum_{d\mid b}\E(\mathcal H,X/N(d),FN(d);\xi,U)\\
 &\ll_{I,\varepsilon}
 D^{\varepsilon/2}(XF)^2\|U\|_{C^J(I_*)}^2.
 \end{aligned}
\]
Here $N(b)\ll_I D$, so the divisor factors are absorbed in $D^{\varepsilon/4}$.
We use the following smooth-weight principle, stated and proved in
Lemma~\ref{lem:smooth-mean-square}: a mean-square bound valid for every
common test function, with a $C^J$ norm, also bounds the corresponding
quadratic sum with a kernel depending on the row, at a cost given by
the kernel's $C^{2J+4}$ norm.
For each fixed $b$, apply Lemma~\ref{lem:smooth-mean-square}
with row index $(f,k)$, coefficient $\mu(f)\ind_{(f,b)=1}$,
and kernel $\mathcal K_{b,f,k}$.
The larger row range $0<N(k)\le\mathcal H$ adds only terms
whose original kernel vanishes. With $q=2J+4$, this gives
\begin{equation}\label{eq:weighted}
 \begin{aligned}
 |\mathcal S_{\xi;B,F}|
 &\ll_{I,\varepsilon}
 D^{\varepsilon/2}\|W\|_{C^q(I)}^2
 \sum_{B\le N(b)<2B}^*\frac{HN(b)}{D^2}(XF)^2\\
 &\ll D^{\varepsilon/2}\|W\|_{C^q(I)}^2
 \frac{HB}{D^2}\,B\,(XF)^2
 =D^{\varepsilon/2}\|W\|_{C^q(I)}^2H,
 \end{aligned}
\end{equation}
where $XF=D/B$ and ideal counting gives $O(B)$ choices of $b$.
Summing the $O((\log D)^2)$ nonempty dyadic ranges bounds
$\mathcal S_\xi$ as required. Sum over the fixed finite set of $\xi$,
include $Z$, and restore the factor $D$ to obtain \eqref{eq:ms}.
The required derivative order depends only on $\varepsilon$.

\end{proof}

\begin{remark}\label{rem:formal-dual}
The formal proof organizes the dual mean squares slightly differently. The excluded primes are
not removed by Lemma~\ref{lem:remove-exclusions}. Instead, an exclusion mask $m$ is carried
inside the dual mean square, through the twist $\chi_n(m^6)=\ind_{(n,m)=1}$. Test functions are
written as profiles $V(\log(\NK(n)/X))$ with compact support and expanded once and for all in
the height twists $\NK(n)^{\mathrm i\theta}$, with bounds polynomial in $\theta$. This is the
expansion behind Lemma~\ref{lem:smooth-mean-square}, built into the definition. Finally, the
rows are weighted by a compactly supported smooth majorant, whose Fourier transform is not
compactly supported but decays rapidly, instead of a weight $\Phi$ with compactly supported
$\widehat\Phi$; the resulting frequency tails are estimated separately. None of these changes affects the exponents.
\end{remark}

